\section{Exact preservation of Shift by eliminating restoration}
\label{sec:qhs-shift}

The immediate parent is PA194QH, whose quantified distribution head has
exchanged premises, whose successful modus ponens leaves its result in the
visible top, and whose Shift tests the cached head. PA194QHS retains these
command semantics. Its new contribution is the exact physical composition
with the previously supplied symmetric comparison method. No individual
component optimization is claimed as new.

Write $T=r_1$, $W=r_2$, $D=r_4$, $H=r_5$, and $X=r_6$.
The resolved QH source has 3,653 instructions and the QHS source has 3,650.
Instructions 0 through 3647 agree literally, including their successors.
The only incoming edge to the removed suffix is the zero branch of
instruction 3644. Instructions 3644--3647 are
\[
 \mathsf{while}\ H>0\ \mathsf{do}\ (H:=H-1;\ T:=T+1;\ X:=X+1).
\]
Consequently $H=0$ at every entry to either suffix. If the values before
this loop were $(T,H,X)=(t,h,x)$, its values on exit are
$(t+h,0,x+h)$. Write $n=x+h$. This statement allows arbitrary incoming
scratch $x$; it does not presume that scratch was zero.

The old suffix restores scratch into the cache and then increments scratch:
\[
 \mathsf{while}\ X>0\ \mathsf{do}\ (X:=X-1;\ H:=H+1);
 \quad X:=X+1;\quad \mathsf{goto}\ 0.
\]
It reaches the common comparator with $(H,X)=(n,1)$. The new suffix is
\[
 H:=H+1;\quad \mathsf{goto}\ 0,
\]
and reaches that comparator with $(H,X)=(1,n)$. The values of $T,W,D$
and every other register agree. The unchanged predecessor loop strictly
decreases $H$; the removed restoration loop strictly decreases $X$.
Their termination and endpoint equations follow by induction on those
natural counters.

\subsection{The comparator is symmetric in its result}

The common source at instructions 0--10 is exactly the following resolved
program. The targets of a decrement are its zero and positive targets.
\begin{center}
\begin{tabular}{rll}
Instruction&Operation&Targets\\\hline
0&decrement $H$&$4,1$\\
1&decrement $X$&$3,2$\\
2&jump&$0$\\
3&jump&$5$\\
4&decrement $X$&$10,5$\\
5&decrement $H$&$7,6$\\
6&jump&$5$\\
7&decrement $X$&$9,8$\\
8&jump&$5$\\
9&jump&$11$\\
10&halt&\\
\end{tabular}
\end{center}

Each completed pass through 0, 1, and 2 reduces both positive counters by
one. If both become zero together, instructions 0 and 4 select the halt
at 10. If one becomes zero first, the other remains positive and control
reaches 5. The loops at 5 and 7 clear both counters and reach 11. Thus,
for all natural $u,v$, the comparator halts exactly when $u=v$; otherwise
it reaches instruction 11 with $(H,X)=(0,0)$. It preserves every other
register. This proves its symmetry as a partial register transformation,
even though the two directions need not have the same duration.

Applying this result to $(u,v)=(n,1)$ and $(1,n)$ proves that both suffixes
halt if $n=1$, and otherwise reach precisely the same source configuration
at 11. The argument covers every suffix-entry store with $H=0$, including
all stores admitted by the predecessor loop.

\subsection{An exact source cost statement}

Count each resolved primitive instruction once, including the final halt
when it is taken. The removed suffix costs $3n+3$ instructions; the new
suffix costs two. The comparator cost from $(u,v)$ is
\[
 C(u,v)=\begin{cases}
 3u+3,&u=v,\\
 2u+v+4,&u>v,\\
 3v+2,&u<v.
 \end{cases}
\]
These equations follow by counting the matched passes and the remaining
clearing loops in the displayed program. Accordingly the old execution
from suffix entry to halt or instruction 11 exceeds the new execution by
\[
 (3n+3+C(n,1))-(2+C(1,n))
 =\begin{cases}
 0,&n=0,\\
 4,&n=1,\\
 2n+4,&n\ge2.
 \end{cases}
\]
The source saving is unbounded. This is a statement about primitive
register instructions. The additional physical fold changes the costs of
other source instructions, so no blanket improvement in literal Turing
machine running time is inferred from it.

\subsection{Command and finite-trace consequences}

At clean command boundaries the unchanged prefix sets scratch to zero
and clears the old top before copying the cached value. The resulting
Shift is therefore, in both QH and QHS,
\[
 S(t,w,d,h)=\begin{cases}
 \mathsf{halt},&h=1,\\
 (h,\operatorname{tl}(w),d,\operatorname{hd}(w)),&h\ne1.
 \end{cases}
\]
The axiom command and modus ponens are unchanged. In particular successful
modus ponens still returns $(b,w,b,0)$ when
$b=d+x$ and $I(t,b)=h$; failure returns $(0,w,b,0)$.

Every visit to the changed suffix satisfies $H=0$, and every such visit
is replaced by a finite execution with the same exit or halt behavior.
Induction on visits, or equivalently on finite execution segments cut at
suffix entry and common instruction 11, proves equality of halting
behavior for the initialized source programs. Neither the shorter source
nor its guard can add a new halt on an unequal comparison. Conversely,
a halt in either source has a finite corresponding execution in the other.
All counter recurrences and finite execution translations are primitive
recursive; their correctness is provable by ordinary induction in PA.

The literal-tape certificates bind the two source endpoints to their exact
machines. The QHS builder reconstructs the pinned QH source and table,
checks all eight simultaneous folds on the final QHS table, and checks
all 4,750 primitive outcomes. The independent executor derives four generic
list lemmas directly from the QHS rows and checks the same outcomes without
an arithmetic-box oracle. Both complete blank initializations are checked.
The new one reaches its source entry after 57,768 transitions; the pinned
QH initialization takes 57,922. These finite startup figures are not evidence
of infinite nonhalting.

Composing the source preservation with these positive-duration physical
simulations gives the ordinary mathematical theorem
\[
 \mathrm{PA}\vdash
 \operatorname{Halt}(\mathrm{PA194QHS},\mathrm{blank})
 \leftrightarrow
 \operatorname{Halt}(\mathrm{PA194QH},\mathrm{blank}).
\]
The inherited QH syntactic consistency calibration then yields
\[
 \mathrm{PA}\vdash
 \neg\operatorname{Halt}(\mathrm{PA194QHS},\mathrm{blank})
 \leftrightarrow \operatorname{Con}(\mathrm{PA}).
\]
This is an ordinary proof using explicitly identified inherited calibration
and symbolic-executor correctness arguments. It is not presented as a new
end-to-end Lean, Rocq, or iPRA kernel theorem.
